\documentclass[11pt]{article}
\usepackage[a4paper,margin=1in]{geometry}
\usepackage{amsmath,amssymb}
\usepackage[hidelinks]{hyperref}
\usepackage{microtype}

\title{Literature Status: The Isotropic Mean-Norm Question}
\author{Automated mathematics run \texttt{fa\_banach\_001}, lane 11}
\date{9 August 2026}

\begin{document}
\maketitle

\begin{center}
\fbox{\parbox{0.90\textwidth}{\textbf{Status: literature-implied answer
(full answer to the extracted mean-norm question).}  This is an
agent-identified scaling implication from a later theorem, not a new proof by
the run.}}
\end{center}

\section*{Original question}

G.~Chasapis, A.~Giannopoulos, and N.~Skarmogiannis ask in
\cite{CGS2019}, arXiv:1906.03719, PDF page 3, whether for every volume-one
isotropic centrally symmetric convex body $K\subset\mathbb R^n$ one has
\[
  L_K\sqrt n\,M(K)\le c(\log n)^b
\]
for some absolute constants $c,b>0$.  Here
\[
  M(K)=\int_{S^{n-1}}\|\theta\|_K\,d\sigma(\theta),
\]
and volume-one isotropy means that the uniform random vector $X_K$ on $K$
satisfies
\[
  \operatorname{Cov}(X_K)=L_K^2\operatorname{Id}.
\]

\section*{Supporting theorem}

P.~Bizeul states the following as Theorem 1.4 of
\cite{Bizeul2026}, arXiv:2607.29458, PDF page 3, attributing the mean-norm
estimate to Bizeul--Klartag together with the newly available dimension-free
bound on the third-moment parameter $\kappa_n$:
\begin{quote}
If $K_0$ is an isotropic convex body in the probabilistic convention
$\operatorname{Cov}(X_{K_0})=\operatorname{Id}$, then
\[
  M(K_0)\le C\sqrt{\frac{\log n}{n}},
\]
where $C$ is universal.
\end{quote}

\section*{Identification}

Given a body $K$ in the convention of the original paper, set
\[
  K_0=L_K^{-1}K.
\]
Then $X_{K_0}=L_K^{-1}X_K$ in distribution, so
$\operatorname{Cov}(X_{K_0})=\operatorname{Id}$.  For every $a>0$, gauge
homogeneity gives $\|x\|_{aK}=a^{-1}\|x\|_K$.  Consequently,
\[
  M(K_0)=L_K M(K).
\]
Applying the supporting theorem to $K_0$ and multiplying by $\sqrt n$ yields
\[
  \boxed{L_K\sqrt n\,M(K)\le C\sqrt{\log n}.}
\]
Thus the extracted question has an affirmative answer with $b=1/2$.  The
supporting paper records sharpness via the cube, so the logarithmic exponent
$1/2$ is optimal up to universal constants.

\section*{Provenance and scope}

The supporting paper does not cite the 2019 question or say that it is
answering arXiv:1906.03719.  The match is therefore an agent-identified direct
reformulation, classified as a \emph{literature-implied answer}.  This
packet resolves only the explicit mean-norm question on source PDF page 3;
it makes no claim about every broader problem discussed in the source paper.

The search used the exact expression $L_K\sqrt n M(K)$ and the phrases
\emph{isotropic mean norm}, \emph{polylogarithmic}, and \emph{$M(K)$}, then
checked arXiv:2510.20511, 2607.24164, and 2607.29458.  The last paper provides
the decisive theorem in its final unconditional form.

\begin{thebibliography}{9}
\bibitem{CGS2019}
G.~Chasapis, A.~Giannopoulos, and N.~Skarmogiannis,
\emph{Norms of weighted sums of log-concave random vectors},
arXiv:1906.03719 (2019).

\bibitem{Bizeul2026}
P.~Bizeul,
\emph{Optimal $MM^*$ bounds for convex bodies},
arXiv:2607.29458 (2026), Theorem 1.4.
\end{thebibliography}

\end{document}
